% !TeX root = main.tex
\changecolours{ScoutPurple}
\section[Desconstrução e Junção: unwind e lookup]{Desconstrução de Arrays e Junções entre Coleções}

% ------------------------------------------------------------------------------
\begin{frame}[fragile]{O Estágio \texttt{\$unwind}: Desconstrução de Arrays}
  \vspace{0.1cm}
  \begin{columns}[t,onlytextwidth]
    \begin{column}{0.38\textwidth}
      \alert{Mecanismo de Desdobramento}\par
      \footnotesize
      \begin{itemize}
        \item Desmembra um atributo do tipo array, gerando um documento BSON independente para cada elemento contido.
        \item Viabiliza agrupamentos analíticos sobre itens aninhados (ex.: faturamento por SKU de produto).
        \item \textbf{Atenção à Cardinalidade:} Se $10.000$ pedidos possuem em média 3 itens, o \lstinline[language=mongosh, basicstyle=\footnotesize\ttfamily]|$unwind| gera $30.000$ documentos em RAM!
        \item \textbf{Preservação de Vazios:} Por padrão descarta arrays vazios; use \lstinline[language=mongosh, basicstyle=\scriptsize\ttfamily]|preserveNullAndEmptyArrays: true| para retê-los.
      \end{itemize}
    \end{column}
    \begin{column}{0.60\textwidth}
      \alert{Exemplo no E-commerce ShopGrids}\par
      \vspace{0.05cm}
      \begin{lstlisting}[language=mongosh, basicstyle=\scriptsize\ttfamily]
db.pedidos.aggregate([
  { $match: { status: "entregue" } },
  { $unwind: "$itens" },
  {
    $group: {
      _id: "$itens.sku",
      total_unidades: { $sum: "$itens.quantidade" },
      receita_gerada: { $sum: "$itens.subtotal" }
    }
  },
  { $sort: { receita_gerada: -1 } },
  { $limit: 5 }
])
      \end{lstlisting}
    \end{column}
  \end{columns}
\end{frame}

% ------------------------------------------------------------------------------
\begin{frame}[fragile]{O Estágio \texttt{\$lookup}: Junção Externa (\textit{Left Outer Join})}
  \vspace{0.1cm}
  \begin{columns}[t,onlytextwidth]
    \begin{column}{0.48\textwidth}
      \alert{Cruzamento entre Coleções}\par
      \footnotesize
      \begin{itemize}
        \item Executa junção externa (\textit{Left Outer Join}) entre coleções do mesmo banco de dados.
        \item Os registros relacionados são acoplados na forma de um \textbf{array} no documento resultante.
        \item Caso não haja correspondência, o array resultante retorna vazio (\lstinline[language=mongosh, basicstyle=\footnotesize\ttfamily]|[]|).
        \item \textbf{Exigência Crítica de Índice:} O campo \lstinline[language=mongosh, basicstyle=\footnotesize\ttfamily]|foreignField| na coleção de destino \textbf{exige índice B-Tree}, evitando varredura global (\texttt{COLLSCAN}) a cada documento!
      \end{itemize}
    \end{column}
    \begin{column}{0.48\textwidth}
      \alert{Sintaxe de Junção no MongoSH}\par
      \vspace{0.05cm}
      \begin{lstlisting}[language=mongosh, basicstyle=\scriptsize\ttfamily]
db.pedidos.aggregate([
  { $match: { status: "entregue" } },
  {
    $lookup: {
      from: "clientes",
      localField: "cliente_id",
      foreignField: "_id",
      as: "cliente"
    }
  },
  { $unwind: "$cliente" },
  {
    $project: {
      numero_pedido: 1,
      "cliente.nome": 1,
      valor_total: 1
    }
  }
])
      \end{lstlisting}
    \end{column}
  \end{columns}
\end{frame}
